\lecture{3.9}{Equivalence and Extensions}
\label{lec:3-9}

\subsection{Equivalence of G-Computation and IPTW Identification}

G-computation and IPTW look different, but they are algebraically equivalent as identification formulas. For $a=1$,
\[
\E\left\{\frac{\ind(A=1)Y}{\Prob(A=1\given W)}\right\}
=
\E\{\E(Y\given A=1,W)\}.
\]
The derivation follows from repeated use of the tower rule:
\[
\E\left\{\frac{\ind(A=1)Y}{g_1(W)}\right\}
=
\E\left[
\E\left\{
\frac{\ind(A=1)Y}{g_1(W)}
\given A,W
\right\}
\right]
\]
\[
=
\E\left[
\frac{\ind(A=1)}{g_1(W)}
\E(Y\given A,W)
\right]
=
\E\left[
\frac{\ind(A=1)}{g_1(W)}
\E(Y\given A=1,W)
\right]
\]
\[
=
\E\left[
\E\left\{
\frac{\ind(A=1)}{g_1(W)}
\E(Y\given A=1,W)
\given W
\right\}
\right]
\]
\[
=
\E\left[
\frac{\E(Y\given A=1,W)}{g_1(W)}
\E\{\ind(A=1)\given W\}
\right]
=
\E\{\E(Y\given A=1,W)\}.
\]

\subsection{Conditioning on the Propensity Score}

The propensity score is a balancing score. Under conditional randomization and positivity,
\[
(\Yone,\Yzero)\ci A\given g(W),
\]
where $g(W)=\Prob(A=1\given W)$.

The intuition is that within groups of people who have the same probability of treatment, the covariate information relevant to treatment assignment has been balanced sufficiently for exchangeability.

The ATE can therefore also be identified by
\[
\ATE
=
\E\left[
\E(Y\given A=1,g(W))
-
\E(Y\given A=0,g(W))
\right],
\]
where the outer expectation is over the population distribution of $g(W)$.

\subsection{Collider Bias}

Colliders are a major reason that "adjust for everything" is not a valid causal strategy. Consider a picnic example:
\begin{itemize}
  \item $A=1$ means chicken salad, $A=0$ means egg salad;
  \item $Y=1$ means a person already has influenza;
  \item $C=1$ means fever after the picnic.
\end{itemize}
Suppose egg salad is tainted and causes fever, influenza also causes fever, and the sandwich does not cause influenza. Then $C$ is a collider:
\[
A\to C\leftarrow Y.
\]

\begin{figure}[H]
\centering
\begin{tikzpicture}[node distance=1.8cm,>=Latex]
  \node[draw, circle, fill=green!5] (A) {$A$\\sandwich};
  \node[draw, rectangle, fill=gray!12, below=.55cm of A] (a) {$a$};
  \node[draw, circle, fill=blue!5, right=of A] (C) {$C(a)$\\fever};
  \node[draw, circle, fill=orange!8, right=of C] (Y) {$Y$\\flu};
  \draw[->, thick] (a) -- (C);
  \draw[->, thick] (Y) -- (C);
\end{tikzpicture}
\caption{Conditioning on fever opens a path between sandwich type and influenza.}
\end{figure}

Since $Y$ is not downstream of $A$, there is no causal effect of sandwich on influenza. Unconditionally, $A$ and $Y$ may be independent. But conditioning on fever induces dependence: among people with fever, knowing the person ate chicken salad makes influenza more likely because chicken salad did not explain the fever.

\begin{keyidea}
Conditioning can create bias when the conditioned variable is a collider or a descendant of a collider.
\end{keyidea}

\subsection{Causal Effect among Compliers}

In some randomized studies, assignment $Z$ is randomized but treatment actually taken $A$ is not fully controlled. Some participants comply with assignment, while others do not. A common causal target is the effect among compliers:
\[
\E\{Y(1)-Y(0)\given A(1)=1,A(0)=0\},
\]
where $A(z)$ is the treatment that would be taken under assignment $Z=z$.

This estimand is not an ordinary ATE in the full population. It is an effect in a latent subgroup defined by joint potential treatment behaviors under two assignment levels.

\subsection{Principal Stratification}

Principal stratification defines subgroups using potential values of an intermediate variable. For example, in a vaccine study, one might ask for the effect of vaccination on growth among children who would have been infected without vaccination:
\[
\E\{G(1)-G(0)\given I(0)=1\}.
\]
Here $I(0)$ is infection status under no vaccination and $G(a)$ is growth under vaccination level $a$.

Principal strata can be scientifically appealing, but they often involve cross-world or partially latent quantities. Identification typically requires assumptions beyond standard exchangeability and positivity.

\subsection{Transition from Identification to Estimation}

This chapter has produced observed-data parameters for causal quantities. For example,
\[
\E\{\Ya{a}\}
=
\E\{\E(Y\given A=a,W)\}
\]
and
\[
\E\{\Ya{a}\}
=
\E\left\{
\frac{\ind(A=a)Y}{\Prob(A=a\given W)}
\right\}.
\]
These formulas do not yet tell us how to estimate the quantities in finite samples. They tell us what must be estimated: outcome regressions, propensity scores, covariate distributions, and expectations.

\begin{estimationbox}
The next step is estimation. Once identification has converted counterfactual parameters into observed-data parameters, we can construct estimators such as plug-in G-computation estimators, IPTW estimators, and later doubly robust estimators.
\end{estimationbox}
